\section{Related Work}

\subsection{Clinical Gait Analysis and Pose-Based Representations}

Clinical gait analysis supports assessment of motor impairment, functional severity, and rehabilitation outcomes~\cite{perry1992gait,baker2006gait,simon2004gait}. Video-based pose estimation offers an accessible alternative to laboratory motion capture by converting ordinary recordings into joint trajectories~\cite{stenum2021applications,stenum2024clinical}. Markerless toolkits and smartphone-based pipelines further broaden access to quantitative movement analysis~\cite{mathis2018deeplabcut,uhlrich2023opencap}. Validation studies have compared video-derived temporal and kinematic measures with instrumented walkways or motion-capture systems across different gait patterns and clinical cohorts~\cite{verlekar2019validation,yamamoto2021openpose,sabo2022validity}. These studies also show that viewpoint, occlusion, pose-estimation error, and population differences remain important sources of measurement uncertainty.

Gait resources differ in population, acquisition setup, skeleton format, and label definition~\cite{kidzinski2020video,khalil2022ataxia,jun2020pathological}. Pose data provide a shared semantic interface, but do not make these resources clinically equivalent. GaitEval therefore preserves each resource's native representation and dataset-defined task.

\subsection{Visual and Pose-Based Pathological Gait Analysis}

Pose-based gait systems use biomechanical features or learned temporal representations for supervised prediction. Graph models such as ST-GCN, 2s-AGCN, and MS-G3D are widely used for skeleton analysis~\cite{yan2018stgcn,shi2019twostream,liu2020disentangling}, but are generally tied to fixed datasets, skeletons, and label spaces.

Early low-cost RGB and RGB-D systems measured or classified pathological gait using semantic gait features, silhouettes, or tracked skeletons~\cite{dolatabadi2017pathological,verlekar2018automatic,ortells2018vision}. Parkinsonian assessment has since ranged from single-camera kinematic measurement and natural-gait monitoring to learned skeleton--silhouette models for clinical severity estimation~\cite{shin2021quantitative,sabo2020parkinsonian,zeng2023video}. Related work has also estimated cerebral-palsy gait function and kinematics and quantified gait after stroke from ordinary videos~\cite{kidzinski2020video,pantzar2018cerebral,lonini2022stroke}.

For ataxia, Auto-Gait analyzes clinically recorded gait-task videos from multiple sites, while pose-based work has also evaluated performed ataxic gait~\cite{rahman2023autogait,khalil2022ataxia}. A separate clinical study used video keypoints to distinguish Parkinson's disease from spinocerebellar degeneration, showing that cross-condition visual analysis predates GaitEval~\cite{eguchi2024differentiating}. Broader abnormal-gait resources include simulated pathological patterns and clinically annotated abnormal-gait videos~\cite{jun2020pathological,ranjan2024gavd}. Collectively, these studies establish that visual and pose features can support condition-specific prediction; results from performed or simulated gait, however, do not establish patient-level diagnostic validity.

Cross-dataset gait analysis is also established. Transfer-learning and spatiotemporal classifiers have been evaluated across separately collected pathological-gait datasets~\cite{verlekar2018transfer,albuquerque2021spatiotemporal}. Auto-Gait includes leave-one-site-out evaluation, and CARE-PD formalizes within-dataset, cross-dataset, leave-one-dataset-out, and adaptation protocols~\cite{rahman2023autogait,adeli2025carepd}. GaitEval therefore does not claim domain-shift difficulty as a new finding. Instead, it tests whether general-purpose LLMs use supplied pose evidence under same-input controls. Logistic regression and a linear support vector machine receive the same deterministic features, and the ATAXIA--PATHOGAIT transfer experiment is limited to that compatible dataset pair.

\subsection{Language Models for Human Motion and Clinical Gait}

Motion-language systems such as MotionGPT, MotionLLM, and MoChat connect language models with motion through learned tokenization or cross-modal alignment~\cite{jiang2023motiongpt,chen2024motionllm,mo2024mochat}. Gait-specific models have also been developed for biometric recognition and Parkinsonian scoring~\cite{yang2025gaitllm,wang2025agir}. BiomechGPT and BioGait-VLM further connect language-aligned representations with clinically relevant movement analysis~\cite{yang2025biomechgpt,chen2026biogait}. These methods rely on task-specific motion-language training rather than testing off-the-shelf LLM evidence use.

The closest concurrent work evaluates general-purpose LLMs using text-encoded kinematic waveforms~\cite{dindorf2026gaitllm}. More broadly, prior work has examined LLMs for structured and tabular inputs~\cite{fang2024llms} and multimodal instruction following~\cite{liu2023llava,yin2024survey}. GaitEval instead evaluates heterogeneous gait resources through controlled pose-derived views, same-feature linear references, and adaptation experiments. Its purpose is to separate limited motion signal from representation difficulty and failed LLM evidence use, while restricting transfer conclusions to the tested dataset pair.
